\cotefile{19}
\scanpage{36}
% L1, upper left
\uline{\unsure{XI|XII}} 6.2. \unsure{----|\ldots}

% L2–3
\[
T^i(X)\unsure{\mathord{-}|\to}T'{}^i(X)\unsure{,|;}
\quad \deleted{\text{\unsure{dire pourquoi est} \illegible{1}}}
\quad\begin{aligned}&\text{\unsure{iso} pour }i\le n\\
&\text{mono pour }i=n+1\end{aligned}
\]
\margin{\unsure{!?|?}}

% L4
\[
\unsure{0}\to A-B-C\to 0
\]

% L5–6
\diagram{\begin{tikzcd}[ampersand replacement=\&,column sep=small,row sep=small]
\&T^{i-1}(B)\arrow[r]\arrow[d]\&T^{i-1}(C)\arrow[r,dash]\arrow[d]\&T^i(A)\arrow[r,dash]\arrow[d]\&T^i(B)\arrow[r]\arrow[d]\&T^i(C)\arrow[d]\\
\text{(I)}\&T'{}^{i-1}(B)\arrow[r,dash]\&T'{}^{i-1}(C)\arrow[r,dash]\&T'{}^i(A)\arrow[r,dash]\&T'{}^i(B)\arrow[r,dash]\&T'{}^i(C)
\end{tikzcd}}

% L7–12, a)
a) Supposons qu'on sache que $\varphi^i$ un isom pour $i\le n-1$ \inserted{mono pour $i=n$.} Et qu'il y a un syst. de cog. pour lesquels \deleted{\unsure{$i$}\illegible{1}} $\varphi^n$ est un isom. Alors le diagramme (I) pour $i=n$ montre que \deleted{\unsure{pour}} $\varphi^n$ est un isom pour \uline{tous} les \deleted{\unsure{g}} arguments. Plus \unsure{préc.|gén.}, si $A$ est \inserted{et qu'il $\exists\ A\to B$ mono, et que c'est vrai pour $\varphi^n(B)$, c'est vrai pour $\varphi^n(A)$.}

% L13–17, b)
b) Supposons qu'on sache que $\varphi^i$ est un isom pour $i\le n$, et $\varphi^{n+1}$ un mono pour un syst. de cog. Alors le diagramme (I) avec $i=n+1$ montre que $\varphi^{n+1}$ est un mono. pour \uline{tous} les \uline{arguments}. Plus \unsure{préc.|gén.}, \inserted{si $\exists\ A-B$ mono, et si c'est vrai pour $B$, c'est vrai pour $A$.}

% L18–24, a')
a') $\overline{\mathrm M}$ hypothèse préliminaire sur les $\varphi^i$ que en a) (isom si $i\le n-1$, mono si $i=n$). On veut savoir si $\varphi^n(B)$ est un \uline{isom} (i.e. un \unsure{épim|section}) pour $B$ donné. Si on a une suite exacte $0\to A\to B\to C\to 0$, il suffit de le vérifier pour $\varphi^n(A)$ et $\varphi^n(C)$ ? Oui, \uline{pourvu que} $\varphi^{n+1}(A)$ soit injectif, \uuline{non} en général.

% L25–27, b')
b') $\overline{\mathrm M}$ hyp \unsure{préliminaire} que b) ($\varphi^i$ un isom si $i\le n$). On veut savoir si $\varphi^{n+1}(B)$ est un mono. Suffit-il de le vérifier pour $\varphi^{n+1}(A)$ et $\varphi^{n+1}(C)$ ? \uuline{Oui}.

% archive sheet number: 35
\scanpage{37}
% L1–2
\deleted{Enfin on conclut que tt morphisme projectif [simple] est $\in M$ \unsure{grâce} à \unsure{e)}.}

% L3–4
\deleted{\uline{Corollaire} Supposons \deleted{de plus} \inserted{au lieu de c)} la condition suivante satisfaite :}

% L5–11
\deleted{c') Soit \unsure{$g$|$\varnothing$} : $X'\xrightarrow{f}X\xrightarrow{g}Y$ \inserted{des} morphismes \deleted{\unsure{de préschémas}} propres, $X_0$ un sous-préschéma fermé de $X$, $U=X-X_0$, $U'=f^{-1}(U)$. \deleted{et} Supposons que $U'\unsure{\mathord{-}|\to}U$ soit un isomorphisme, \deleted{et} que $X_0\xrightarrow{g_0}Y$ \deleted{soit} et $gf:X'-Y$ soient $\in M$. Alors $g\in M$.}

% L12–13
\deleted{Alors tt morphisme propre $f:X-Y$ \deleted{est} $\in M$.}

% L14–15
\deleted{En effet, il est trivial que c') implique \unsure{a)|c)}, donc tt morphisme projectif est $\in M$.}

% L16–22
\deleted{On obtient le cas d'un morphisme \deleted{\unsure{quelcon|quelconque}} propre quelconque en utilisant \unsure{e)|c)}, \unsure{[}qui permet \deleted{\unsure{de|du}} supposer $Y$ affine,\unsure{]} \deleted{et c'} et d'utiliser \inserted{la récurrence} \deleted{le lemme de} \deleted{Chow, puis c'} noethérienne. Le lemme de Chow donne \deleted{\unsure{donc}} joint à c') donne alors le résultat voulu.}

% archive sheet number: 36
\scanpage{38}
% L1–4, c)
c) Supposons les $\varphi^n$ des isom \deleted{\unsure{pour $i\le n$|fonctoriels}}, \deleted{\unsure{un épimorphisme|un épimorphisme surjectif}} Quand peut-on conclure que $\varphi^{n+1}$ est un mono ? Réponse : il suffit que \deleted{\unsure{$\varphi$}} $T^{n+1}$ soit effaçable. Donc si les $T^i$ \unsure{son|sont} effaçables pour $i>i_0$, et les $\varphi^i$ des \unsure{d° \upshape isom|isom} pour $i$ \deleted{\unsure{petit}} $<i_0$, un

% L5–7, crossed block
\deleted{d) Supposons les $\varphi^n$ des isom, pour \unsure{$i\le n$}, un mono, en \unsure{dim $n+1$}, Quand peut-on affirmer que $\varphi^{n+1}$ est \unsure{$\overline{\mathrm M}$} un isom ? \deleted{Quand} \inserted{Il suffit que} $T'{}^{n+1}$ \deleted{\unsure{$(B)$}} soit coeffaçable}

% L8–10, resumption of c)
mono pour $i\unsure{>|\ge}i_0$, et $\varphi^i$ un épim pour $i\le n$ (où $n\ge i_0$) alors $\varphi^i$ est un \deleted{\unsure{épim|épi}} \uline{isom} pour $i\le n$, un \unsure{épim|épi} en dim $n+1$.

% L11–13
Revenons sur a), b). On a des propriétés \deleted{\unsure{$I$}} $I_n$, $J_n$ d'un argument $A\in\mathrm{Ob}\,C$,
\[
\left(\begin{aligned}
 I_n(A)&:\varphi^{\deleted{n}\inserted{i}}(A)\ \text{isom si }i\le n\\
 J_n(A)&:\varphi^i(A)\ \text{isom si }i\le n,\ \text{mono}\\
 &\hspace{7em}\text{si }i=n+1
\end{aligned}\right]
\]

% L14–16
\[
\left\{\begin{aligned}
 J_n&\Rightarrow I_n\Rightarrow J_{n-1}\quad\ldots\\
 I_n\ \text{pour tt }A,\ +J_n\ \text{pour cog.}&\to J_n\ \text{pour tout } A\\
 J_n\ \text{pour tt }A,\ +I_{n+1}\ \text{pour tt cog}&\Rightarrow I_{n+1}\ \text{pour tt }A
\end{aligned}\right.
\]

% L17–20
\uline{\unsure{Concl|Concl.}} Pour vérifier $I_n$ (resp $J_n$) pour tt $A$, sachant que c'est vrai pour $n$ petits, il suffit de vérifier $I_m(B)$ (resp $J_m(B)$) pour $m\le n$ et tt cog. $B$.

% L21, double rule
\rule{3cm}{0.4pt}\\[-1.5ex]\rule{3cm}{0.4pt}

% L22–27
On avait ici des arguments où on ``\unsure{montait|remontait}'' les degrés. Dans les questions de dim \unsure{ch.|coh.}, au contraire, il faut ``descendre''. La question est de donner des critères pour qu'un $T^i$ soit nul (pour tt \uline{argument}). Ici on doit utiliser les faits suivants (\unsure{tenant déjà|tenant d'abord|savoir aussi dans} au fait que $T^i$ est

% archive sheet number: 37
\scanpage{39}
% L1–2
\deleted{\deleted{\uline{Lemme \unsure{2|1}}} \inserted{\uline{Proposition}} Soit \deleted{\unsure{$F$|$J$}} $M$ un ensemble de morphismes \deleted{(\unsure{simples projectifs})} de préschémas loc. noethériens. On suppose que}

% L7–8, a); circled addition with pointer
\deleted{a) La projection $f:P^1_Y\to Y$ \deleted{est} \inserted{et $\mathrm{id}_Y:Y\to Y$ sont} $\in M$.}

% L9–10
\deleted{b) Si \deleted{\unsure{$f$}} $X\xrightarrow{f}Y\xrightarrow{g}Z$ $f,g\in M$, alors \deleted{\unsure{on est de}} $gf\in M$}

% L11–15
\deleted{c) Si $X=P^r_Y\xrightarrow{g}Y$ est le morphisme structural canonique, et \deleted{si} $X'=P^1_Y\times P^{r-1}_Y$ est muni \deleted{du} \deleted{\unsure{$X'$}} morphisme birationnel bien connu $X'\overset{f}{\unsure{\mathord{-}|\to}}X$, et si $gf\in M$, alors $g\in M$.}

% L5–7, inner relocation arrow to just above Sous ces conditions
\inserted{\deleted{d) Si $X\xrightarrow{f}Y\xrightarrow{g}Z$ \deleted{sont} tels que $f$ soit une immersion fermée, $g\deleted{f}\in M$ [et $gf$ simple] alors $gf\in M$}}

% L3–4, outer relocation arrow to just above Sous ces conditions
\inserted{\deleted{\unsure{e)|a)} Pour un morphisme $f:X\to Y$, l'appartenance à $M$ est locale sur $Y$}}

% L16–17
\deleted{Sous ces conditions, tout morphisme [simple] projectif $f:X\to Y$ est $\in M$.}

% L18–22
\deleted{En effet, on voit par récurrence sur $r$, utilisant a), b), c), que les $P^{\deleted{\unsure{n}}r}_Y\xrightarrow{t}Y$ sont $\in M$. Utilisant d), il s'ensuit que pour tt sous-préschéma fermé $X$ de $P^n_Y$ [simple$/Y$] la projection \deleted{sur} $Y$ est $\in M$.}

% archive sheet number: 38
\scanpage{40}
% L1–3
semi-exact :
\[
\left\{\begin{aligned}
 &\text{Si }T^i(B)=0,\ \text{alors }T^i(A)=0\ \text{pour tt facteur direct }A\text{ de }B\\
 &\text{Si } A\text{ a une filtration finie à facteurs }A_\alpha\\
 &\qquad\text{tels que }T^{\unsure{i|v}}(A_\alpha)=0\ \text{pour tt }\alpha,\ T^i(A)=0
\end{aligned}\right.
\]

% L4–12
Pour les foncteurs cohomologiques de faisceaux de $l$-torsion \uline{constructibles} sur $X$ (\unsure{pas qct pas séparé}) ces remarques \deleted{déjà} impliquent déjà que pour avoir $T^i(F)=0$ pour tt tel $F$, il suffit de le prouver pour un $F$ de la forme $i_!(p_*(\unsure{\mathbb Z|Z}/l\mathbb Z)_{Y'})$\inserted{$=q_*j_!(\mathbb Z/l\mathbb Z)_{Y'})$}, où $Y$ est un sous-préschéma constructible de $X$, $Y'$ un rev. étale fini de $Y$ de degré constant, et $p:Y'\unsure{\mathord{-}|\to}Y$ et $i:Y\to X$ les morphismes can. On \unsure{pt|vt} prendre $j$ une immers. ouverte dans $X'$ fini sur $X$.

% L12–14, left margin
\diagram{\begin{tikzcd}[ampersand replacement=\&,row sep=small,column sep=small]
Y'\arrow[r,"j"]\arrow[d,"p"']\&X'\arrow[d,"q"]\\
Y\arrow[r,"i"']\&X
\end{tikzcd}}

% L13–21
Toute sous-catégorie pleine de la catégorie des faisceaux de $l$-torsion constructibles sur $X$, stable par facteurs directs, par extension, et contenant les faisceaux $q_*(j_!(\mathbb Z/l\mathbb Z)_{Y'})$, précédents, est toute // Si par ailleurs cette catégorie est \deleted{\unsure{aussi}} telle que pour tte suite exacte $0-A-B-C-0$, si $B$ et $C$ y sont, alors $A$ y est --- p. ex si c'est une sous-catégorie \uline{épaisse} --- alors il suffit de vérifier que les $q_*(\mathbb Z/l\mathbb Z)_{X'}$ y sont, $X'$ fini sur $X$.

% L22–24
Si $X$ \unsure{noeth}, on peut même \unsure{tj se|tjrs} prendre $X'$ intègre. Si $X$ \deleted{\unsure{est de caractéris} \illegible{1}} \inserted{de plus} univ. japonais, on peut prendre $X'$ intègre normalisé.

% archive sheet number: 39
\scanpage{41}
% L1–5
\deleted{\deleted{\unsure{sur|sont}} en un pt \unsure{au} dessus de $s$, i.e. de $X_0$. Donc pour voir que c'est un isom \unsure{en} pt \unsure{gén.|géom.} $\unsure{z|\bar z|Z}$ de $X_0$, il suffit de voir que \deleted{\unsure{c'est}} \inserted{en} lui appliquant $\unsure{\mathbb R|R}j^!$, où $j:X_0\to Y$, on trouve un isomorphisme en $z$. Or on a}

% L6–7
\diagram{\begin{tikzcd}[ampersand replacement=\&,column sep=large,row sep=small]
\deleted{\mathbb Rj^!(\mathbb Ri^!(\mu_{\alpha}))}\arrow[r]\arrow[d,phantom,"\deleted{\unsure{\wr\!\vert|\Vert|\wr}}" description]\&\deleted{\mathbb Rj^!(\unsure{\Lambda|A}_Y[-2]]}\arrow[d,phantom,"\deleted{\Vert}" description]\\
\deleted{\mathbb Rh^!\deleted{\unsure{i^!}}(\unsure{\mu_X|N_X|\mu_\alpha})}\&{}
\end{tikzcd}}

% L7, left
\diagram{\begin{tikzcd}[ampersand replacement=\&,column sep=small]
\deleted{X_0}\arrow[r,"\deleted{j}"]\arrow[rr,bend right=23,dash,"\deleted{h}"']\&\deleted{Y}\arrow[r,"\deleted{i}"]\&\deleted{X}
\end{tikzcd}}

% archive sheet number: 40
\scanpage{42}
% L1, upper left
$\overline{\underline{\mathrm{XII}}}$

% L1–4
\uline{Remarque} \inserted{Dans 6.5.} Les hypothèses de quasi-compacité et quasi-séparation ne servent que pour \deleted{\unsure{faire}} prouver l'équivalence de la condition a) avec les conditions en apparence plus faibles b), b') etc.

% L5–10
L'équivalence de ces conditions entre elles, et les implications \deleted{de} a) \deleted{\unsure{implique}} vers b) \unsure{énoncées|s'énonce} dans 6.5. bis, restent valables sans cette hypothèse et méritent d'être énoncées une bonne fois \ldots. Dans 6.6. également, on n'a \unsure{d'ailleurs} besoin de la quasi-compacité et quasi-séparation \unsure{;|,}

% L11–15
de plus, la catégorie de tous les faisceaux d'ens. (resp. de Groupes de torsion \unsure{$\mathscr L$|$\mathcal L$} etc) peut ad libitum s'y remplacer par telle ou telle sous-catégorie --- manifestement on n'a pas \unsure{donné|d'ennui} l'énoncé \unsure{bouffe-tout} abstrait.

% L15–23
Dans le lemme 6.7. de descente également, on n'utilise pas la quasi-compacité et quasi-séparation, (sauf, \uline{soit-disant}, pour passer à la limite \deleted{\unsure{dans la}} \inserted{à la catégorie} des faisceaux de $\unsure{\mathbb Z|Z}/n\mathbb Z$-Modules, mais c'est sans doute inutile \deleted{\unsure{quitte à en être affirmé}} \inserted{--- et par l'argument à la main, utilisant \uline{\unsure{VIII|VII|XIV}} 5.2.,} \deleted{\unsure{--- ce qui c'est que cette vérification}} pour déduire que l'hom de changement de base en dim $n+1$ est encore injectif --- mais cela doit aussi sortir directement.

% archive sheet number: 41
\scanpage{43}
% L1–3
groupes sur $\bar X$, également ind-fini en vertu de (IX 1.6. (ii)). \handwritten{\'}{E}crivons maintenant l'homomorphisme naturel de suites exactes d'ensembles pointés, où l\deleted{\unsure{es}}\handwritten{a} première flèche horizontale de chaque ligne est injective:\handwritten{\tikz[baseline=-.3ex]{\draw (0,.04) circle (.038) (0,.078)--(.03,.29);}}
\margin{\handwritten{\tikz[baseline=-.3ex]{\draw (0,.04) circle (.038) (0,.078)--(.03,.29);} réf}}

% L4–5
\diagram{\begin{tikzcd}[ampersand replacement=\&,column sep=small,row sep=small]
\deleted{\unsure{0}}\&H^1(X,\handwritten{\pi_*}(\bar F))\handwritten{\arrow[r]\arrow[d]}\&H^1(\bar X,\bar F)\handwritten{\arrow[r]\arrow[d]}\&H^0(X,R^1\handwritten{\pi_*}(\bar F))\handwritten{\arrow[d]}\\
\deleted{\unsure{0}}\&H^1(Y,\handwritten{\pi}_{Y\handwritten{*}}(\bar F_{\bar Y}))\handwritten{\arrow[r]}\&H^1(\bar Y,\bar F_{\handwritten{\bar Y}})\handwritten{\arrow[r]}\&H^0(Y,R^1\handwritten{\pi}_{Y\handwritten{*}}(\bar F_{\bar Y})),
\end{tikzcd}}

% L6–13
cf (\handwritten{\unsure{3?2.|3.2.}}). Compte tenu de l'hypothèse sur les homomorphismes de changement de base faites dans (i) et (ii), et de la co\deleted{\unsure{nditi}}\handwritten{nclusi}on de (i), la pre\handwritten{mi}ère flèche verticale s'identifie à la flèche de restriction $H^1(X,F)\handwritten{\to}H^1(\deleted{\unsure{X}}Y,F_Y)$ dont nous voulons établir la surjectivité, et la dernière flèche verticale est \handwritten{in}jective\handwritten{;} d'autre part, d'après l'hypothèse (ii), la deuxième flèche verticale est bijective. On en conclut aussitôt la \deleted{\unsure{bijectiv}} surjectivité de la première flèche verticale, ce qui achève de prouver (ii). \deleted{\unsure{En}}

% L14–17
(iii) \deleted{(\unsure{Esquisse|quisse} abrégée). On adopte pour abréger les notations des catégories dérivées (sinon, il faut passer à des résolutions injectives..). On considère l'homomorphisme canonique $F\handwritten{\to}R\handwritten{\pi_*}(\bar F)$, où $\bar F=\handwritten{\pi^*}(F)$, et son mapping-cylinder $G$, de sorte qu'on a un triangle exact}

% L18, triangle
\diagram{\begin{tikzcd}[ampersand replacement=\&,column sep=tiny,row sep=small]
\&\deleted{G}\handwritten{\arrow[dl]}\&\\
\deleted{F}\handwritten{\arrow[rr]}\&\&\deleted{R\handwritten{\pi_*}(\bar F)}\handwritten{\arrow[ul]}{\qquad\deleted{\text{.}}}
\end{tikzcd}}

% L19–23
\deleted{Pour \unsure{pouver|ppouver} \deleted{\unsure{le théorême}} l'assertion pour $F$, on est ramené à la prouver pour $R\handwritten{\pi_*}(\bar F)$ et $G$ séparément, grâce à la suite exacte de cohomologie et au lemme des cinq\deleted{\unsure{s}}. Pour $G$, on note que $G$ est \deleted{à cohomologie} \deleted{\unsure{faisceaux de}} nul\deleted{\unsure{le}} sur $U$, et pour un tel complexe de faiceaux, la suite spectrale de cohomologie nous ramène au cas d'un faisceau concentré sur $Y$, qui est trivial.}

% L24–25
\deleted{Le cas de $R\handwritten{\pi_*}(\bar F)$ se traite en notant que \deleted{\unsure{$\underline H R(\bar F)\quad\underline H$}} l'homomorphisme $\underline H^\bullet(X,R\handwritten{\pi_*}(\bar F))\handwritten{\to}H^\bullet(Y,h^{\handwritten{*}}(R\handwritten{\pi_*}(\bar F))$ s'identifie, grâce à l'hypothèse sur l'homo}

% L26–28, bottom in pen
\inserted{\handwritten{Adapter la démonstration de \uline{XII} 8.3. (iii); on se ramène si on veut au cas $\pi$ \uline{surjectif}, en remplaçant $\bar X$ par $\bar X\amalg Y$ !}}

% archive sheet number: 42
\scanpage{44}
% L1–2, upper right
Foncteurs \unsure{localisants|localisant} et Catégories de fractions

% L3, faint pencil to the left
4 p + 5 p

% archive sheet number: 43
\scanpage{45}
% L1, 3)
3)
\diagram{\begin{tikzcd}[ampersand replacement=\&,column sep=large]
\underline A\arrow[r,shift left,"T"]\&\underline B\arrow[l,shift left,"S"]\&\operatorname{Hom}(x,Sy)\overset{\sim}{\unsure{\mathord{-}|\to}}\operatorname{Hom}(Tx,y)
\end{tikzcd}}
\margin{$\left\{\vphantom{\rule{0pt}{9em}}\right.$}

% L2–3
$x\in\mathrm{Ob}\,\underline A$ conditions équivalentes \uline{si} $S$ \uline{pleinement fidèle}

% L4
\textcircled{a} $x$ est isomorphe à un objet de la forme $Sy$

% L5
\textcircled{b} $x\xrightarrow{\sim}STx$

% L6
\textcircled{c} $\forall\,x'\in A,\ \operatorname{Hom}(x',x)\xrightarrow{\sim}\operatorname{Hom}(Tx',Tx)$

% L7–8
($b\Leftrightarrow c\Rightarrow a$ sans hyp sur $S$, $a\Rightarrow b$ si on suppose $S$ pleinement fidèle).

% L9, 4)
4)
\diagram{\begin{tikzcd}[ampersand replacement=\&,column sep=large]
\underline A\arrow[r,shift left,"T"]\&\underline B\arrow[l,shift left,"S"]\&\operatorname{Hom}(x,Sy)\overset{\sim}{-}\operatorname{Hom}(Tx,y).
\end{tikzcd}}

% L10
Conditions équivalentes

% L11
a) $S$ pleinement fidèle

% L12
b) $TS\overset{\sim}{-}\mathrm{id}_{\underline B}$ isomorphisme

% L13–17
c) $T$ est un foncteur de passage au quotient, i.e. Si on considère \deleted{\unsure{l'ens|l'ens.}} $\unsure{M|\mathcal M}$ des $u\in\mathrm{Fl}\,\underline A$ tels que $T(u)\in\unsure{\mathrm{Flis}|\mathrm{Fliso}|\mathrm{Fl}_{\rm iso}|\mathrm{Isom}}(\underline B)$, et le foncteur $\underline A/M\to\underline B$, ceci est une \uline{équivalence} de catégories.

% archive sheet number: 44
\scanpage{46}
% L1–2, 2)
2) $\underline A\xrightarrow{T}\underline A/\underline C=\underline B$

Conditions équivalentes sur $T$ i.e. sur $\underline C$

% L3–4
\textcircled{a} $T$ admet un foncteur adjoint $S$ :
\[
\operatorname{Hom}(Tx,y)\overset{\sim}{\unsure{\mathord{-}|\to}}\operatorname{Hom}(x,Sy)
\]

% L5–9
\textcircled{b} $\forall\, x\in\underline A$, il existe un plus grand sous-objet $x'$ de $x$ qui est $\in\underline C$, \uline{et} si ce dernier est nul, il existe un monomorphisme de $x$ dans un objet $\underline C$-fermé

% L10–12
\textcircled{c} $\forall\,x\in\mathrm{Ob}\,\underline A$, il existe un hom
\[
x\longrightarrow\tilde x
\]
qui est un $\underline C$-isom, avec $\tilde x$ $\underline C$-fermé.

% L13–14, addition between the two lists
\inserted{Alors le hom en c) est unique à isom unique près, c'est aussi $x\longrightarrow STx=Lx$}

% L15–16, 1)
1) $\underline A\overset{T}{-}\underline A/\underline C=\underline B$, $x\in\mathrm{Ob}\,\underline A$. Conditions équivalentes sur $x$ (pour $\underline C$ i.e. $T$ \unsure{donné|donnés|donnée})

% L17
\textcircled{a} $\forall\,x'\in A,\ \operatorname{Hom}(x',x)\xrightarrow{\sim}\operatorname{Hom}(Tx',Tx)$

% L18–19
\textcircled{b} $\forall$ $x'-x''$ un $\underline C$-isom, on a
\[
\operatorname{Hom}(x'',x)\xrightarrow{\sim}\operatorname{Hom}(x',x)
\]

% L18–20, lower left
% slanted
\margin{\unsure{i.e.} pour tt $\underline C$-isom \unsure{$u$} \illegible{1} \unsure{$x''$}, est un isom \uline{surjectif}, \uline{\unsure{non}}, \unsure{et|i.e.} \deleted{\unsure{est}} \unsure{$\exists$} un inverse à droite}

% L20–21, torn lower edge
\textcircled{c} Tout sous objet de $x$ qui est $\in\underline C$ est nul, et \inserted{pour} tt \unsure{$y\supset x$|$y>x$|$x\subset x'$} \unsure{tel que} $\unsure{y/x|x'/x}\in\underline C$ \deleted{\unsure{est} \illegible{1}} $x$ est facteur direct !

% archive sheet number: 45
\scanpage{47}
% L1–4, printed heading
% letterhead
\begin{center}
INSTITUT\\
DES HAUTES ÉTUDES SCIENTIFIQUES\\
5, ROND-POINT BUGEAUD, PARIS XVI\textsuperscript{e}\\
KLÉBER\quad 48-56
\end{center}

% L5–8
$C$ catégorie abélienne avec $\underline{\lim}$ exactes et \unsure{gén.}

$C''\subset C'$ sous-catégories abéliennes localisantes, i.e.

$1^\circ$) épaisses

$2^\circ$) stables par $\underline{\lim}$

% L9–10
Foncteurs
\[
\begin{aligned}
\underline{\Gamma}'(F)&=\text{plus grand sous-objet de }F\quad\in\mathrm{Ob}\,C'\\
\underline{\Gamma}''(F)&=\rule{6.5em}{.3pt}\qquad\qquad\in\mathrm{Ob}\,C''
\end{aligned}
\]

% L9–12, circled note with leaders to both functor signs
\margin{définis comme adjoints de $i'_*,i''_*$ noter $i^{\prime !},i^{\prime\prime !}$}

% L11
\[
\underline{\Gamma}'/\underline{\Gamma}''(F)=\underline{\Gamma}'(F)/\underline{\Gamma}''(F)
\]

% L12
Foncteurs
\diagram{\begin{tikzcd}[ampersand replacement=\&]
 C' \arrow[r,shift left,"T"] \& C''/C' \arrow[l,shift left,"S"]
\end{tikzcd}}
et foncteur section

% L12–15, second marginal note, linked to the circled note
% slanted
\margin{distinguer dans la notation avec $\underline{\Gamma}',\underline{\Gamma}''$ foncteurs à val. \uline{sur $C$}}

% L13–15
d'où foncteur
\[
\Lambda=ST:C'\longrightarrow C'
\]
avec
\[
\mathrm{id}_{C'}\longrightarrow\deleted{\unsure{ST}}\Lambda
\]

% L16–20
caractérisé par

$1^\circ$) $x'\unsure{\mathord{-}|\to}\Lambda(x')$ est un $C''$-isom

$2^\circ$) $\Lambda(x')$ est ``$C''$ fermé'', i.e. pour tt hom $y'-z'$ dans $C'$, $\operatorname{Hom}_{C'/C''}(z',\Lambda x')-\operatorname{Hom}_{C'/C''}(y',\Lambda x')$ est bijectif.

% L21–23
\uline{Lemme 1} $\underline{\Gamma}',\underline{\Gamma}'',\Lambda$ exacts à gauche, $\Gamma'/\Gamma''$ transforme

\deleted{\unsure{on a} $\Gamma'(F)\longrightarrow\Lambda\Gamma'(F)$}

\unsure{mono. en mono.}

% L24–25
Pour $\Lambda$, par \unsure{univers...|universalité} (\unsure{trivial ?}). Pour $\Gamma'$ on note que \deleted{\illegible{1}} $\operatorname{Hom}_C(i'_*(x'),x)\simeq\operatorname{Hom}_{C'}(x',\underline{\Gamma}'(x))$, OK.

% archive sheet number: 46
\scanpage{48}
% L1–4
Soit $x\unsure{\mathord{-}|\to}y$ un mono, considérons $x'',y''$ les plus grands sous-objets de $x,y$ $\in\mathrm{Ob}\,C''$. Il faut prouver que
\diagram{\begin{tikzcd}[ampersand replacement=\&,row sep=0.1em,column sep=0.3em]
x \&[-1em] \subset \&[-1em] y\\
\cup \& \& \cup\\
x'' \& \subset \& y''
\end{tikzcd}}
$x''=x\cap y''$, or c'est trivial.

% L5
\uline{Cor} $\underline{\Gamma}'(x)/\underline{\Gamma}''(x)\hookrightarrow R^0(\underline{\Gamma}'/\underline{\Gamma}'')(x)$

% L6–7
Considérons pour $x\in\mathrm{Ob}\,C$ variable
\[
\underline{\Gamma}'x\longrightarrow\Lambda\underline{\Gamma}'x
\]

% L8–9
d'où un hom
\[
\underline{\Gamma}'/\underline{\Gamma}''(x)\longrightarrow\Lambda\underline{\Gamma}'(x)
\]

% L10–12
et comme $\Lambda\underline{\Gamma}'$ est exact à g.
\diagram{\begin{tikzcd}[ampersand replacement=\&]
\underline{\Gamma}'x \arrow[d] \arrow[dr] \& \\
R^0\underline{\Gamma}'/\underline{\Gamma}'' \arrow[r] \& \Lambda\Gamma'
\end{tikzcd}}

% L13–14
\uline{Prop} C'est là un \uline{isom} (?)

équivaut :

% L13–17, left margin
% slanted
\margin{\uline{Faux en} \uline{général} ! \uline{Ok. si les} \uline{objets injectifs} de \uline{$C''$ sont} \uline{injectifs dans $C$}}

% L15–17
\uline{Cor} Pour tt $x\in\mathrm{Ob}\,C$, on a ceci

a) $R^0(\underline{\Gamma}'/\underline{\Gamma}'')(x)$ est $C''$-\uline{fermé} [dans $\underline C'$, ou dans $\underline C$ c'est pareil...]

b) $\underline{\Gamma}'x/\underline{\Gamma}''(x)\longrightarrow R^0(\underline{\Gamma}'/\underline{\Gamma}'')(x)$ est un $\underline C''$-isom.

% L18–19
Cela équivaut aussi (c'est \unsure{animé} par un \unsure{Mons|Mais} \unsure{la|le} \deleted{\unsure{vérif}} \unsure{vérifier})

% L20–21
\uline{Cor} Si $x$ est injectif, alors $\underline{\Gamma}'(x)/\underline{\Gamma}''(x)\xrightarrow{\sim}\Lambda\underline{\Gamma}'(x)$, i.e. $\underline{\Gamma}'(x)/\underline{\Gamma}''(x)$ est \deleted{$\unsure{\underline{\Gamma}'|\underline C'}$} $\underline C''$-fermé.

% L22–23
En effet, $x/\underline{\Gamma}''(x)$ est $C''$-fermé, \deleted{et on voit} \inserted{[en effet, il est injectif par loc. cit \unsure{lem. 2|lex. 2|cor. 2} et cor. 3}

% archive sheet number: 47
\scanpage{49}
% L1, heading
Conditions à envisager

% L2
(i) $\underline C\xrightarrow{i_*}\underline A$ transforme injectifs en injectifs

% L2–7, brace beside (i)–(iii)
% slanted
\margin{\unsure{équivalentes|équivalents} $\left\{\vphantom{\rule{0pt}{4em}}\right.$}

% L3–5
(ii) Pour tout objet injectif $x$ de $\underline A$, le plus grand sous-objet $x'$ de $x$ qui est $\in\mathrm{Ob}\,\underline C$ est encore injectif \uline{dans} $A$.

% L6–8
(iii) Pour tout objet de $\underline C$, son enveloppe injective dans $\underline A$ est dans $\underline C$ [i.e. $\underline C$ stable par enveloppes injectives] -- donc \unsure{néc.|aussi|a. c.} injective dans $\underline C$.

% L8–14, two arrows below the brace, ending at (iv) and (v)
\diagram{\begin{tikzcd}[ampersand replacement=\&,row sep=1.2em,column sep=1em]
{} \arrow[d] \arrow[dd,bend right=22,Rightarrow] \& \\ {} \& \\ {} \&
\end{tikzcd}}

% L9–13
(iv) \deleted{Pour tout $x$} L'hom naturel
\[
R^0(\mathrm{id}/\unsure{i_*|[{}_*}i^!)\longrightarrow ST
\]
est un isomorphisme, i.e. \deleted{pour} $x$ objet injectif de $\underline A$, \deleted{$\unsure{x/i|x/i^{!}|\text{si}}$} si $x'$ \inserted{$=i_*i^{\prime !}(x')$} est le plus grand sous-objet $\in\mathrm{Ob}\,\underline C$, alors $x/x'$ est $C''$-fermé.

% L12–14, margin crossing the longer arrow
% slanted
\margin{cf Gabriel, cor. 3}

% L14
(v) $T:\underline A\unsure{\mathord{-}|\to}\underline A/\underline C$ transforme injectifs en injectifs.

% L15–17
$(\mathrm{i})\Rightarrow(\mathrm{ii})$ Car \deleted{$x'=$} pour \uline{tt} $x$, $x'=i_*i^!(x)$, or $i^!(x)$ [étant adjoint de $i_*$ qui est exact] transforme injectifs en injectifs

% L18–22
$(\mathrm{ii})\Rightarrow(\mathrm{i})$ Car si $x\in\underline C$, soit $\bar x$ une enveloppe injective de $x$ dans $\underline A$, \deleted{et} considérons \deleted{$\unsure{i_*|i}$}$\bigl(i^!(\bar x)\bigr)$, c'est un objet \uline{injectif} dans $\underline C$ contenant $x$, et une extension essentielle du dit [car elle l'est dans $C$] donc c'est une enveloppe injective de $x$ \uline{dans} $\underline C$.

% L23–24
Si $\bar x$ est injectif, il est \unsure{en} égal à $i^!(\bar x)$. Si (ii) est vérifié, $i_*i^!(\bar x)$ est injectif, i.e. $i_*(x)$ injectif.

% archive sheet number: 48
